
\newpage
\section{Reading(5.1)} 
\subsection{The GPU Computing Era.}
The article describes the historical evolution of GPU computing, starting with a demand for real time graphic rendering in the late 1990s. It highlights how, with the addition of CUDA additional flexibility was added, enabling an expansion into a multitude of additional computing domain. The authors then give a quick introduction into the CUDA programming model, which shall not be explained here as it corresponds to the lecture contents. Then follows a description of several aspects of the Fermi architecture, talking about PCIe as the GPU-CPU interconnect, the memory hierarchy, streaming multiprocessors, parallel thread exectuion and the concept of CUDA cores. It concludes with a pledge for using GPU CPU coprocessing, showing how parallel intensive as well as mostly sequential programs can leverage the advantages of having both components available.\\
\newline
What can be taken away from the article is, that has been shown since, further evolution with regards to more general/generic use of GPUs might enable even more domains to leverage the advantages of GPU computing. Additionally, with the GPU CPU coprocessing, the advantages CPUs have will still remain valid, even with the further evolution of GPUs. The key takeaway is to try and make use of the availability of both components to achieve the best performance possible.\\
\newline
In general, we agree with the notion of the article to harness the power of GPU CPU coprocessing. An aspect though which is neglected in the article is the additonal space, power and monetary requirements of having a system equipped with both components. The final projection with regards to a 50\% increase in compute power per year might have been correct at the time of writing, but is very optimistic for the current state of hardware.